\documentclass[11pt]{article}
\usepackage[T1]{fontenc}
\usepackage[utf8]{inputenc}
\usepackage{lmodern}
\usepackage[margin=1in]{geometry}
\usepackage{microtype}
\usepackage{longtable}
\usepackage{amsmath,amssymb}
\usepackage[colorlinks=true,urlcolor=blue]{hyperref}
\title{Mathematical Manuscripts\\\large Overview and Verification Guide}
\author{Richard J. Reyes\\Independent Researcher}
\date{October 26, 2026}
\begin{document}
\maketitle
\begin{abstract}
Mathematical manuscripts contributing to an ongoing Wave Confinement Theory
(WCT) research program. Formal verification is used for rigorous development
and correction, not abandonment of foundational physical objectives.
Analytic arguments, Lean declarations, numerical evidence and physical
interpretations remain distinct. External peer review is not implied.
\end{abstract}
\section*{Research purpose and interpretation}
The objective is to investigate whether confinement, resonance, curvature,
phase dynamics and nonlinear fields can support a mathematically consistent,
experimentally testable physical theory. Independent mathematical results
do not themselves validate WCT physics; obstructions only exclude their
precisely quantified constructions. Failed routes should lead to explicit
new mathematical or empirical tests, not assumed repairs or a compulsory
engineering-only pivot. Consistent ghost-free parent models and viable
relativistic structure remain separate research obligations.
See \url{https://github.com/rickyjreyes/math/blob/main/RESEARCH_DIRECTION.md}.
\section*{Flagship: fourth-order cubic--quintic variational PDE}
On $H^2(\mathbb R^2;\mathbb C)$, for $k,\beta,\gamma>0$, consider
\[
H(u)=\|(\Delta+k^2)u\|_2^2-\frac{\beta}{2}\|u\|_4^4+
\frac{\gamma}{3}\|u\|_6^6.
\]
The primary family presents sufficient conditions for strict subadditivity,
minimizer existence and orbital stability in the complete minimizer set.
These claims concern a declared independent mathematical Hamiltonian.
\section*{Result families}
\begin{longtable}{p{0.12\textwidth}p{0.78\textwidth}}
\textbf{Family/ID}&\textbf{Contents}\\\hline
F01&Strict binding, minimizers, orbital stability and negative trials\\
F02&Higher-derivative constraint counting and Sobolev UV obstruction\\
F03&Spectral--phase exact identity and closed-curve counterexample\\
F04&Nonlinear finite Fourier subspace non-invariance\\
F05&Positive curvature whole-space nonattainment\\
F06&Finite-wavenumber dispersion minimum\\
F07&Exact weighted phase-locking variational result\\
F08&Rank and signature constraints on scalar-gradient tensors\\
F09&11-by-31 structural local identifiability obstruction\\
F10&No $L^2$ state supported on a single Fourier shell\\
F11&E13/E14 band-pass dissipation; E16 linear spectral selection; E59 fixed-annulus projection; E62 Fourier fraction\\
E61&Centered continuum Swift--Hohenberg spectral threshold\\
M6B&Nonuniqueness counterexamples for curvature closure and fixed-mass inverse map\\
F12&E26 conditional historical additive-denominator $L^2$ estimate\\
WCT E68&Localized-energy counterexample and corrected $L^n$, $H^2$ estimates\\
\end{longtable}
\paragraph{Regularized curvature scope.}
E26 treats the historical additive-denominator expression and assumes a
uniform positive lower bound on its modulus. The separately catalogued
M2/E17 result treats a corrected modulus-squared, phase-covariant
regularization; the two operators should not be conflated.
\section*{Independent reproduction}
Browse the catalogue and PDFs at
\url{https://github.com/rickyjreyes/math/blob/main/CONTENTS.md}.
Lean declarations and explicit verification status:
\url{https://github.com/rickyjreyes/math/blob/main/lean/formalization.yaml}.
The maintained Lean project includes formalizations of F01 nonlinear-PDE manuscript results. Exact scope and the remaining analytic-only manuscripts are recorded in the manifest; mathematical proof alone does not validate physical interpretation.
\end{document}
